\documentclass[10pt,a4paper]{amsart}
\usepackage[margin=19mm]{geometry}
\usepackage{amsmath,amssymb,mathtools}
\usepackage[scr=rsfs,cal=euler]{mathalfa}
\usepackage{xfrac}
\usepackage[unicode,hidelinks,bookmarks=false]{hyperref}
\newcommand{\ie}{i.e.\ }
\newcommand{\cf}{cf.\ }
\newcommand{\resp}{respectively\ }
\newcommand{\Subsection}{\subsection}
\providecommand{\nobreakdash}{\nobreak\hbox{-}}
\setlength{\emergencystretch}{2em}
\multlinegap0pt
\usepackage{amssymb}
\newtheorem{lemma}{Lemma}
\title{Thick Difference Sets of Haar Null Compact Sets in Locally Compact Groups}
\author{Chuck Akemann}
\date{Source: arXiv:2603.27479v1 (CC BY 4.0)}
\begin{document}
\maketitle

\noindent\emph{Source and licence.} Thick Difference Sets of Haar Null Compact Sets in Locally Compact Groups. Version arXiv:2603.27479v1 (CC BY 4.0), distributed by arXiv under the Creative Commons Attribution 4.0 International licence, http://creativecommons.org/licenses/by/4.0/, as recorded in the arXiv licence metadata for this exact version and shown on the abstract page https://arxiv.org/abs/2603.27479v1. Full licence notice: \texttt{licenses/arxiv-2603.27479-CC-BY-4.0.txt}.\par\smallskip

\noindent\textit{Editorial scope.} Counted unit: the article's lemma lemma (source0.tex lines 133--138). The article's complete original proof of it is delivered with it (source0.tex lines 140--229, one \texttt{proof} environment of 90 lines). No mathematics is written, repaired or re-derived: the statement and the proof are the article's own lines, in the article's own notation, and no retained line is substituted or re-flowed.  No further source-local context is needed: every cross-reference of every retained line resolves inside the two delivered spans.  Nothing was omitted from any retained span, so the package carries no omission record.  The retained lines cite 1 works, whose entries are transcribed verbatim from the article's own bibliography source in the e-print and delivered with short editorial labels recorded in the item's reference mapping.  No retained line contains a figure environment, an image-inclusion command or drawing code, so no image is delivered and the package declares no asset dependency.  Editorial formatting only, in addition: an AMS \texttt{amsart} preamble that restates the article's own declarations and macros used by the retained lines, namely the environments \texttt{lemma} on separate counters, together with the standard packages those lines need.  The retained environments print in this document as Lemma 1 in order of appearance, whereas the article prints the same items with the numbering of its own full document.  The complete native source of the article as distributed in the arXiv e-print for this exact version is bundled unchanged as \texttt{sources/197352/source0.tex}.
\begin{lemma}
Let \(G\) be an infinite compact totally disconnected group.
Let \(O\) be a neighborhood of the identity.
Then there exists a compact set \(K \subset O\) with \(m(K)=0\)
such that \(KK^{-1}\) contains a neighborhood of the identity.
\end{lemma}

\begin{proof}
Since \(G\) is compact and totally disconnected,
it is profinite and admits a neighborhood basis at the identity
consisting of open normal subgroups (see, e.g., Hofmann and Morris \cite[Theorem~1.34]{HM}).

We construct a descending sequence of open normal subgroups
\[
N_1 \supset N_2 \supset \cdots,\quad N_1 \subset O
\]
inductively so that the successive quotients \(Q_k := N_k/N_{k+1}\) satisfy
\[
\frac{A}{\sqrt{|Q_k|}} \le \frac12
\]
(where \(A=4/\sqrt{3}\)). This is possible because \(G\) is infinite profinite:
given \(N_k\), the open normal subgroups contained in \(N_k\) still form a neighborhood basis of the identity, so we can skip intermediate subgroups to ensure the index \([N_k:N_{k+1}]=|Q_k|\) is arbitrarily large.

Let \(H_k = N_1/N_{k+1}\) and let \(\Phi_k : N_1 \to H_k\) be the canonical quotient map.
Let \(\pi_k : H_{k+1} \to H_k\) be the natural projection between quotients, so that \(\Phi_k = \pi_k \circ \Phi_{k+1}\).

We construct subsets \(S_k \subset H_k\) recursively so that
\[
S_k S_k^{-1} = H_k, \quad \pi_k(S_{k+1}) = S_k,
\]
and writing
\[
\eta_k := \frac{|S_k|}{|H_k|},
\]
we have
\[
\eta_{k+1} \le \eta_k \frac{A}{\sqrt{|Q_{k+1}|}} \le \frac{\eta_k}{2}.
\]

First choose \(S_1 \subset H_1\) such that
\[
S_1 S_1^{-1} = H_1 \quad\text{and}\quad |S_1| \le A \sqrt{|H_1|}.
\]

Suppose \(S_k\) has been constructed. Choose
\(T_{k+1} \subset Q_{k+1}\) such that
\[
T_{k+1} T_{k+1}^{-1} = Q_{k+1} \quad\text{and}\quad |T_{k+1}| \le A \sqrt{|Q_{k+1}|}.
\]

For each \(x \in S_k\), choose a lift \(\tilde{x} \in H_{k+1}\)
with \(\pi_k(\tilde{x}) = x\), and let \(\tilde{S}_k\) denote
the set of these lifts. Define
\[
S_{k+1} := \tilde{S}_k T_{k+1}.
\]

Because \(T_{k+1} T_{k+1}^{-1} = Q_{k+1}\) and \(Q_{k+1}\) is normal in \(H_{k+1}\), we have
\[
S_{k+1} S_{k+1}^{-1} = \tilde{S}_k T_{k+1} T_{k+1}^{-1} \tilde{S}_k^{-1} = \tilde{S}_k Q_{k+1} \tilde{S}_k^{-1} = \tilde{S}_k \tilde{S}_k^{-1} Q_{k+1}.
\]
Since \(\pi_k(\tilde{S}_k \tilde{S}_k^{-1}) = S_k S_k^{-1} = H_k\), the product covers all of \(H_k\), and multiplying by the kernel \(Q_{k+1}\) yields all of \(H_{k+1}\). Thus \(S_{k+1} S_{k+1}^{-1} = H_{k+1}\).

Moreover,
\[
\eta_{k+1} = \frac{|S_{k+1}|}{|H_{k+1}|} \le \frac{|S_k||T_{k+1}|}{|H_k||Q_{k+1}|} \le \eta_k \frac{A}{\sqrt{|Q_{k+1}|}} \le \frac{\eta_k}{2}.
\]

Thus \(\eta_k \le \eta_1 / 2^{k-1} \to 0\).

Define
\[
K_k := \Phi_k^{-1}(S_k) \subset N_1.
\]
Because \(\Phi_k = \pi_k \circ \Phi_{k+1}\) and \(\pi_k(S_{k+1}) = S_k\), it follows that \(K_{k+1} \subset K_k\). Each \(K_k\) is compact, and
\[
m(K_k) = \eta_k\, m(N_1) \to 0.
\]

Let
\[
K = \bigcap_{k=1}^\infty K_k.
\]
Then \(K\) is compact, \(K \subset N_1 \subset O\), and \(m(K)=0\).

Let \(g \in N_1\). For each \(k\), since \(H_k = S_k S_k^{-1}\),
there exist \(a_k, b_k \in S_k\) such that
\[
\Phi_k(g) = a_k b_k^{-1}.
\]
Thus \(K_k \cap gK_k \neq \emptyset\) for all \(k\).

Since these sets are nested and compact, their intersection is nonempty,
so \(g \in KK^{-1}\).

Hence \(N_1 \subset KK^{-1}\).
\end{proof}

\begin{thebibliography}{99}
\bibitem[HM]{HM} K. H. Hofmann and S. A. Morris, \textit{The Structure of Compact Groups}, De Gruyter, 2013.
\end{thebibliography}
\end{document}
